\documentclass[tikz,border=8pt]{standalone}
\usepackage{ctex}
\usepackage{amsmath}
\usepackage{pgfplots}
\pgfplotsset{compat=1.18}
\usetikzlibrary{calc,positioning,arrows.meta}

%% 图 opt-p2-05-2：步长选择对比
%% 相同函数三条 GD 轨迹：左太小（缓慢）/ 中合适（平滑）/ 右太大（震荡/发散）

\begin{document}
\begin{tikzpicture}

%% ===== 子图 1：步长太小（缓慢收敛）=====
\begin{scope}[xshift=0cm]

  % 坐标轴
  \draw[->, gray!50, thin] (-2.2, 0) -- (2.2, 0)
    node[right, font=\scriptsize] {$x_1$};
  \draw[->, gray!50, thin] (0, -2.0) -- (0, 2.2)
    node[above, font=\scriptsize] {$x_2$};

  % 椭圆等高线
  \foreach \r in {0.5, 1.0, 1.5, 2.0}{
    \draw[gray!35, thin] (0,0) ellipse ({1.8*\r} and {\r});
  }

  % 极小值点
  \fill[red!80!black] (0,0) circle (2.5pt);
  \node[font=\scriptsize, below right] at (0,0) {$\mathbf{x}^*$};

  % 太小步长的轨迹（缓慢，密集点）
  \coordinate (A0) at (1.8, 1.0);
  \coordinate (A1) at (1.55, -0.75);
  \coordinate (A2) at (1.32, 0.58);
  \coordinate (A3) at (1.12, -0.44);
  \coordinate (A4) at (0.95, 0.34);
  \coordinate (A5) at (0.81, -0.26);
  \coordinate (A6) at (0.69, 0.20);

  \draw[orange!80!black, thick] (A0)--(A1)--(A2)--(A3)--(A4)--(A5)--(A6);
  \draw[orange!80!black, thick, dashed, ->] (A6) -- (0.4, 0.1);

  \foreach \pt in {A0,A1,A2,A3,A4,A5,A6}{
    \fill[orange!80!black] (\pt) circle (1.5pt);
  }
  \fill[orange!80!black] (A0) circle (2.5pt)
    node[above right, font=\scriptsize] {$\mathbf{x}_0$};

  % 标题
  \node[font=\small\bfseries, orange!80!black] at (0, -2.5) {步长太小};
  \node[font=\scriptsize] at (0, -2.9) {$\alpha \ll 1/L$，缓慢收敛};

\end{scope}

%% ===== 子图 2：步长合适（平滑收敛）=====
\begin{scope}[xshift=5.5cm]

  % 坐标轴
  \draw[->, gray!50, thin] (-2.2, 0) -- (2.2, 0)
    node[right, font=\scriptsize] {$x_1$};
  \draw[->, gray!50, thin] (0, -2.0) -- (0, 2.2)
    node[above, font=\scriptsize] {$x_2$};

  % 椭圆等高线
  \foreach \r in {0.5, 1.0, 1.5, 2.0}{
    \draw[gray!35, thin] (0,0) ellipse ({1.8*\r} and {\r});
  }

  % 极小值点
  \fill[red!80!black] (0,0) circle (2.5pt);
  \node[font=\scriptsize, below right] at (0,0) {$\mathbf{x}^*$};

  % 合适步长的轨迹（适度之字，较快收敛）
  \coordinate (B0) at (1.8, 1.0);
  \coordinate (B1) at (0.9, -0.5);
  \coordinate (B2) at (0.45, 0.22);
  \coordinate (B3) at (0.18, -0.08);
  \coordinate (B4) at (0.06, 0.03);

  \draw[green!60!black, thick, ->] (B0)--(B1);
  \draw[green!60!black, thick, ->] (B1)--(B2);
  \draw[green!60!black, thick, ->] (B2)--(B3);
  \draw[green!60!black, thick, ->] (B3)--(B4);
  \draw[green!60!black, thick, dashed, ->] (B4) -- (0.01, 0.005);

  \foreach \pt in {B0,B1,B2,B3,B4}{
    \fill[green!60!black] (\pt) circle (2pt);
  }
  \fill[green!60!black] (B0) circle (2.5pt)
    node[above right, font=\scriptsize] {$\mathbf{x}_0$};

  % 标题
  \node[font=\small\bfseries, green!60!black] at (0, -2.5) {步长合适};
  \node[font=\scriptsize] at (0, -2.9) {$\alpha \approx 1/L$，平滑收敛};

\end{scope}

%% ===== 子图 3：步长太大（震荡/发散）=====
\begin{scope}[xshift=11cm]

  % 坐标轴
  \draw[->, gray!50, thin] (-2.2, 0) -- (2.2, 0)
    node[right, font=\scriptsize] {$x_1$};
  \draw[->, gray!50, thin] (0, -2.0) -- (0, 2.2)
    node[above, font=\scriptsize] {$x_2$};

  % 椭圆等高线
  \foreach \r in {0.5, 1.0, 1.5, 2.0}{
    \draw[gray!35, thin] (0,0) ellipse ({1.8*\r} and {\r});
  }

  % 极小值点
  \fill[red!80!black] (0,0) circle (2.5pt);
  \node[font=\scriptsize, below right] at (0,0) {$\mathbf{x}^*$};

  % 步长过大：震荡且越来越大
  \coordinate (C0) at (1.0, 0.5);
  \coordinate (C1) at (-1.8, -0.9);
  \coordinate (C2) at (1.95, 1.0);
  % 已超出等高线范围，用箭头示意发散

  \draw[red!80!black, very thick, ->] (C0)--(C1);
  \draw[red!80!black, very thick, ->] (C1)--(C2);
  \draw[red!80!black, very thick, dashed, ->] (C2) -- ++(0.3, 0.15)
    node[right, font=\scriptsize, red!80!black] {发散！};

  \fill[red!80!black] (C0) circle (2.5pt)
    node[above right, font=\scriptsize] {$\mathbf{x}_0$};
  \fill[red!80!black] (C1) circle (2pt)
    node[below left, font=\scriptsize] {$\mathbf{x}_1$};
  \fill[red!80!black] (C2) circle (2pt)
    node[above right, font=\scriptsize] {$\mathbf{x}_2$};

  % 感叹号警示
  \node[font=\large\bfseries, red!80!black] at (-1.0, 1.5) {!};

  % 标题
  \node[font=\small\bfseries, red!80!black] at (0, -2.5) {步长太大};
  \node[font=\scriptsize] at (0, -2.9) {$\alpha > 2/L$，震荡 / 发散};

\end{scope}

%% ===== 整体标题 =====
\node[font=\large\bfseries] at (8.0, 3.0)
  {步长 $\alpha$ 选择对梯度下降收敛的影响};

%% ===== 底部说明 =====
\node[draw=gray!50, thin, fill=gray!5, rounded corners=3pt,
      font=\scriptsize, inner sep=5pt, align=center]
  at (8.0, -3.6)
  {稳定条件：$\alpha \leq 1/L$（$L$ 为梯度 Lipschitz 常数）\quad
   最优步长：$\alpha = 2/(\mu+L)$（强凸时）};

\end{tikzpicture}
\end{document}
